% Limitations section, written 2026-08-23, held out of the paper on request
% (page budget). Insert immediately before \section{Conclusion}. It depends on
% \ifFragileNone / \ifFragilePattern and on \label{sec:related}.

\section{Limitations}
\label{sec:limitations}
The evidence comes from one testbed. Five datacentres, one workload trace, one wind dataset and one
forecaster produce every number reported here. The mechanism argument is stated for any scheduler
whose exogenous prediction enters through an identifiable channel, and nothing in the method is
specific to wind or to carbon, but the demonstration is not.

Training is open-book on a single wind window, and that window sits at the tenth percentile of the
year's green supply. The policy therefore learns from one realisation of the exogenous process it is
asked to reason about. Evaluation uses three windows the training never saw, spanning low, middle and
high scarcity, which tests whether the learned behaviour transfers. It does not widen what the
behaviour was learned from.

Corruption is injected rather than observed. Blend and Shuffle are synthetic degradations applied at
evaluation, chosen because each has a plausible operational cause, and Shuffle in particular
reproduces a misconfigured site-to-feed mapping. No trace of a real forecaster degrading in
production drives any result here, so the magnitude of the penalty is a property of the injected
corruption as much as of the policy.

The containment advantage holds under deterministic decoding and closes under sampling. Deterministic
decoding is the deployment form, for reasons of capacity planning and service certification that hold
independently of this work, so the boundary is placed where deployment places it. That is an argument
about operational practice rather than a demonstration that the mechanism survives a sampled policy,
and a scheduler deployed with sampled actions would not inherit the result.

\ifFragileNone\else\ifFragilePattern\else
One seed produced a policy whose greedy decode over-commits to deferral while its sampled decode
completes the workload. A single instance locates nothing, and it is reported here rather than
analysed.
\fi\fi

The training-time decomposition and the deployment-time auditor share a principle and not an
implementation. The auditor computes no gradients, changes no weights, and would function in front of
a policy trained by any method. Presenting them together is a claim about one idea applied at two
time scales, and a reader who wants the decomposition evaluated alone should read
Section~\ref{sec:main} without Appendix~\ref{app:auditor}.

Hindsight credit assignment is the nearest prior treatment of exogenous influence on the return, and
it is discussed in Section~\ref{sec:related} but not run as a baseline. The comparison set here
covers risk-reshaping objectives, which target the same symptom through the objective rather than
through credit, and a policy trained without the forecast at all.

Finally, the environment is a simulator. Carbon, completion and utilisation are computed from a model
of the fleet, and no result has been reproduced on physical hardware.

